\chapter{Neyman's confidence constructions and comparisons}
\label{sec:p15-neyman-confidence-constructions}

Neyman's 1937 examples separate three questions: whether a procedure
has the prescribed coverage, whether its realized intervals are short,
and whether it seldom includes false parameter values
\cite[pp.~356--378]{Neyman1937Confidence}. Here
\(0<\alpha<1\) denotes the source's \emph{confidence coefficient};
the corresponding testing error probability is \(1-\alpha\).
The acceptance/interval identities are those of
\cref{thm:p15-confidence-test-inversion,thm:p15-neyman-interval-construction}.
All interval endpoints below are included. The particular boundary
completion required for the switched example is defined explicitly.

\section{Example I: two observations from a uniform population}

\begin{hbdefinition}{def:p15-neyman-uniform-systems}{The uniform experiment and three systems}
Let \(X_1,X_2\) be independent uniform variables on \((0,\theta)\),
\(\theta>0\), on sample space \((0,\infty)^2\).
Write
\[
 l=\min(X_1,X_2),\quad L=\max(X_1,X_2),\quad
 S=X_1+X_2,\quad \bar X=S/2,\quad
 \beta=1-\sqrt\alpha,\quad q=\sqrt{1-\alpha}.
\]
For a candidate value \(t>0\), define the acceptance sets
\[
\begin{aligned}
 A_1^{\rm U}(t)&=\{\beta t\leq X_1,X_2\leq t\},\\
 A_2^{\rm U}(t)&=\{qt\leq S\leq(2-q)t\},\\
 A_3^{\rm U}(t)&=\{qt\leq L\leq t\}.
\end{aligned}
\]
Their inverted sets are
\[
\begin{aligned}
 I_1^{\rm U}(x)&=[L,l/\beta]\quad\hbox{if }L\leq l/\beta,
                    \quad\varnothing\hbox{ otherwise},\\
 I_2^{\rm U}(x)&=[S/(2-q),S/q],\\
 I_3^{\rm U}(x)&=[L,L/q].
\end{aligned}
\]
The sum acceptance band has half-width \((1-q)t\) about \(t\).
The sample-space boundary choices in \(A_1^{\rm U}\) and
\(A_3^{\rm U}\) do not change their probability under any
fixed uniform law, and give these exact closed inversions.
\end{hbdefinition}

\begin{hbtheorem}{thm:p15-neyman-uniform-coverage}{proved=annex:p15.neyman.uniform.coverage}{All three coverages and the empty-section failure}
Each displayed system has coverage exactly \(\alpha\) for every
\(\theta>0\). The sum and maximum systems give nonempty intervals
for every sample point. The first system has an empty inverted
set exactly when \(l<\beta L\), an event of probability \(\beta\);
it therefore fails the source's requirement of a nonempty interval
at every sample point, although it is a valid possibly empty
confidence-set procedure.

The sum has density
\[
 f_S(s)=
 \begin{cases}s/\theta^2,&0<s<\theta,\\
 (2\theta-s)/\theta^2,&\theta\leq s<2\theta,\\
 0,&\text{otherwise}.
 \end{cases}
\]
For \(n\geq1\) independent uniform observations and maximum \(L_n\),
\[
 P_\theta(L_n\leq z)=(z/\theta)^n\quad(0\leq z\leq\theta).
\]
Consequently
\[
 I_n^{\rm max}(x)=[L_n,L_n/q_n],\qquad q_n=(1-\alpha)^{1/n},
\]
has coverage exactly \(\alpha\).
\end{hbtheorem}

\begin{hbtheorem}{thm:p15-neyman-uniform-lengths}{proved=annex:p15.neyman.uniform.lengths}{The source's length and repeated-experiment comparisons}
Let \(D_2,D_3\) be the lengths of the sum and maximum intervals.
Then
\[
 D_2=\frac{2(1-q)S}{q(2-q)},\qquad
 D_3=\frac{(1-q)L}{q},\qquad
 \frac{2-q}{4}\leq\frac{D_3}{D_2}<\frac{2-q}{2}<1.
\]
The lower equality occurs when \(X_1=X_2\).
Their expected lengths are
\[
 E_\theta D_2=\frac{2(1-q)\theta}{q(2-q)},\qquad
 E_\theta D_3=\frac{2(1-q)\theta}{3q}.
\]
At \(\alpha=3/4\), the source's observations \((1,1)\)
and \((1/10,19/10)\) both give the sum interval
\([4/3,4]\); their maximum intervals are respectively
\([1,2]\) and \([19/10,19/5]\).
The length ratio is in \([3/8,3/4)\).

In 400 independent ideal uniform experiments with any prescribed
positive scales, each procedure's coverage count is
\(\operatorname{binomial}(400,3/4)\). Thus its probability
of coverage frequency outside \([0.69,0.81]\) is exactly
\[
 \sum_{k=0}^{275}\binom{400}k(3/4)^k(1/4)^{400-k}
 +\sum_{k=325}^{400}\binom{400}k(3/4)^k(1/4)^{400-k}.
\]
The source's particular schedule assigns 155 trials scale 1,
37 scale 2, 8 scale 10, 10 scale 20 and 190 scale 30.
The two counts obtained from the same sample pairs are not independent.
This is a sampling-model calculation, not a report that the source's
random-number-table exercise has been performed.
\end{hbtheorem}

\section{Conditional similarity and Example II}

\begin{hbdefinition}{def:p15-neyman-conditional-family}{The specified conditional-density interface}
Suppose a specified measurable invertible coordinate representation
of the observations is \((T,V)\), with its actual joint density
\[
                    m_{\theta,\eta}(z)\,h_\theta(v\mid z).
\]
Here \(m_{\theta,\eta}\) integrates to 1, and \(h_\theta(\cdot\mid z)\)
is a measurable conditional probability density independent of
the nuisance parameter \(\eta\), for marginally relevant \(z\).
The coordinate representation and any density Jacobian are part of
the hypotheses; algebraic independence alone is not a change-of-variables
theorem.
A measurable family of sections \(B_t(z)\) is conditionally similar
at confidence coefficient \(\alpha\) if
\[
                       \int_{B_t(z)}h_t(v\mid z)\,dv=\alpha
\]
for every relevant \(z\). Its acceptance region consists of observations
whose transformed coordinates satisfy \(V\in B_t(T)\).
\end{hbdefinition}

\begin{hbtheorem}{thm:p15-neyman-conditional-similarity}{proved=annex:p15.neyman.conditional.similarity}{The complete conditional construction of Propositions VIII--IX}
Under the stated factorization, the conditional distribution of \(V\)
given \(T=z\) is independent of \(\eta\). Combining measurable
sections of conditional probability \(\alpha\) produces an acceptance
region of probability exactly \(\alpha\) for every nuisance value.
The sections may be unions of intervals or other measurable sets;
being a single interval is not required.
If the residual coordinate \(V\) has positive Euclidean dimension,
its conditional marginal cumulative function for the first coordinate
is continuous. Selecting an interval of length \(\alpha\) in that
cumulative-probability coordinate constructs such sections, in
infinitely many ways, with joint measurability.
The theorem requires actual sections of the stated conditional
probability. Their existence at arbitrary levels cannot be inferred
for atomic conditional laws merely from sufficiency.
\end{hbtheorem}

\begin{hbdefinition}{def:p15-neyman-triangle-law}{The actual two-parameter triangular density}
Let
\[
 \theta_1>0,\qquad -1<\theta_1^3\theta_2\leq2,\qquad
 \rho=\theta_1^3\theta_2.
\]
On \((0,\infty)^2\) define
\[
 f_{\theta_1,\theta_2}(x_1,x_2)=
 \left(\frac2{\theta_1^2}-\theta_1\theta_2+3\theta_2x_1\right)
                           \mathbf1_{\{x_1+x_2<\theta_1\}}.
\]
Density values on the boundary are set to zero.
This is the distribution of source Example II, with the complete
cubic restriction on the nuisance parameter retained.
\end{hbdefinition}

\begin{hbtheorem}{thm:p15-neyman-triangle-conditional}{proved=annex:p15.neyman.triangle.conditional}{Normalization, the failed radial band and the sufficient coordinate}
The displayed function is a probability density, positive in the
open triangle. For \(0<a<1\),
\[
 P_{\theta_1,\theta_2}\{a\theta_1\leq X_1+X_2\leq\theta_1\}
             =1-a^2+\frac12a^2(1-a)\rho.
\]
It depends on the nuisance value and cannot give one fixed
confidence coefficient for all such values.
The marginal density of \(X_1\) is
\[
 \left(\frac2{\theta_1^2}-\theta_1\theta_2+3\theta_2x_1\right)
                      (\theta_1-x_1)\mathbf1_{\{0<x_1<\theta_1\}},
\]
and, given \(X_1=x_1\), the conditional law of \(X_2\) is
uniform on \((0,\theta_1-x_1)\), independently of \(\theta_2\).
The conditional statement concerns \(0<x_1<\theta_1\);
there is no division by zero at \(x_1=\theta_1\).
\end{hbtheorem}

\begin{hbdefinition}{def:p15-neyman-triangle-intervals}{The fixed-placement family from Example IIa}
For \(0<b\leq1-\alpha\), use
\[
 A_b^{\rm T}(t)=
 \{0<x_1<t:\ b(t-x_1)\leq x_2\leq(b+\alpha)(t-x_1)\}.
\]
It inverts to
\[
             I_b^{\rm T}(x)=
                [x_1+x_2/(b+\alpha),\ x_1+x_2/b].
\]
Let \(I_*^{\rm T}=I_{1-\alpha}^{\rm T}\); explicitly,
\[
             I_*^{\rm T}(x)=[x_1+x_2,\ x_1+x_2/(1-\alpha)].
\]
\end{hbdefinition}

\begin{hbtheorem}{thm:p15-neyman-triangle-coverage}{proved=annex:p15.neyman.triangle.coverage}{Exact nuisance-uniform coverage and the length optimum within the family}
Every \(I_b^{\rm T}\) has coverage \(\alpha\) for every allowed
\((\theta_1,\theta_2)\) and is nonempty at every positive sample.
Its length is
\[
                          \frac{\alpha x_2}{b(b+\alpha)}.
\]
This length strictly decreases as \(b\) increases, so \(b=1-\alpha\)
is uniquely shortest at every sample point within this fixed-placement
family. This assertion specifies its comparison class.
\end{hbtheorem}

\begin{hbdefinition}{def:p15-neyman-switched-fibres}{The switched conditional placement with complete boundary sections}
For candidate \(t>0\), define \(A_{\rm sw}^{\rm T}(t)\) as follows:
\[
\begin{array}{ll}
0<x_1<t/2:&(1-\alpha)(t-x_1)\leq x_2\leq t-x_1,\\
t/2<x_1<t:&0<x_2\leq\alpha(t-x_1),\\
x_1=t/2:&0<x_2\leq t-x_1.
\end{array}
\]
The third line completes the switch boundary. For each fixed
parameter it is a zero-probability line.
For \(u=x_1>0,v=x_2>0\), put
\[
\begin{aligned}
 \ell_{\rm sw}(u,v)&=\max\{u+v,\ u+\min(v/\alpha,u)\},\\
 h_{\rm sw}(u,v)&=\max\{2u,\ u+v/(1-\alpha)\}.
\end{aligned}
\]
\end{hbdefinition}

\begin{hbtheorem}{thm:p15-neyman-switched-intervals}{proved=annex:p15.neyman.switched.intervals}{Every switched case and the source's crossing-length comparison}
The exact inversion of \(A_{\rm sw}^{\rm T}\) is the nonempty
closed interval \([\ell_{\rm sw},h_{\rm sw}]\), for all
\(0<\alpha<1\). It has nuisance-uniform coverage \(\alpha\).
At \(v=\alpha u\), its length is
\[
                       \frac{\max\{2\alpha-1,0\}}{1-\alpha}u,
\]
strictly less than the length \(\alpha^2u/(1-\alpha)\) of
\(I_*^{\rm T}\). At \(v=qu\) with
\[
                      0<q<\frac{\alpha(1-\alpha)}
                                         {1-\alpha+\alpha^2},
\]
its length is \(u-v/\alpha\), strictly greater than that
of \(I_*^{\rm T}\).
At \(\alpha=3/4\), the observations \((1,3/4)\) give
switched interval \([2,4]\) and fixed-placement interval
\([7/4,4]\); \((1,1/10)\) gives respectively
\([17/15,2]\) and \([11/10,7/5]\).
\end{hbtheorem}

\section{The original accuracy and unbiasedness criteria}

\begin{hbdefinition}{def:p15-neyman-shortest-system}{Neyman's shortest-system criterion}
For an exact-coverage interval procedure \(I\), let
\[
                     Q_I(s,\eta;t)=P_{s,\eta}\{t\in I(X)\}.
\]
A shortest system in the source's sense minimizes this probability
for every false target \(t\ne s\), simultaneously, among the
specified class of procedures with confidence coefficient \(\alpha\).
The nuisance values remain quantified over.
This is a false-coverage criterion, not the assertion that every
realized interval or its expected length is smallest.
A one-sided lower estimate \(\ell(X)\) is best in the source's
sense when \(P_{s,\eta}\{\ell(X)\leq t\}\) is minimal for
every \(t<s\), subject to its prescribed coverage; the upper
version minimizes \(P_{s,\eta}\{t\leq u(X)\}\) for \(t>s\).
\end{hbdefinition}

\begin{hbtheorem}{thm:p15-neyman-accuracy-duality}{proved=annex:p15.neyman.accuracy.duality}{The exact reduction to acceptance probabilities}
For \(A_I(t)=\{x:t\in I(x)\}\),
\[
                             Q_I(s,\eta;t)=P_{s,\eta}(A_I(t)).
\]
Thus a simultaneously optimal family of acceptance regions,
whose inversion is an admissible interval procedure, gives a
shortest system. One-sided versions have the same identity with
the corresponding half-line interval.
In addition, when the interval membership graph is measurable,
\[
                             E_{s,\eta}\operatorname{len}(I(X))
                              =\int_{\mathbb R}Q_I(s,\eta;t)\,dt,
\]
allowing infinity. Here \(\operatorname{len}\) denotes Lebesgue length.
Pointwise false-coverage dominance therefore implies expected-length
dominance, while a comparison of realized lengths alone is a
different claim.
\end{hbtheorem}

\begin{hbtheorem}{thm:p15-neyman-uniform-shortest}{proved=annex:p15.neyman.uniform.shortest}{The complete uniform-model shortest-system proof}
The maximum interval \(I_n^{\rm max}\) minimizes false coverage
simultaneously for every true \(s>0\) and candidate \(t>0\), \(t\ne s\),
among all exact-\(\alpha\) procedures. Its coverage function is
\[
 Q(s;t)=
 \begin{cases}
 \alpha(t/s)^n,&0<t\leq s,\\
 1-(1-\alpha)(t/s)^n,&s<t<s/q_n,\\
 0,&t\geq s/q_n.
 \end{cases}
\]
It is strictly below \(\alpha\) at every false target.
The same minimum holds even when the competing acceptance
rules are randomized measurable functions with expectation
\(\alpha\) at the candidate law.
\end{hbtheorem}

\begin{hbtheorem}{thm:p15-neyman-triangle-shortest}{proved=annex:p15.neyman.triangle.shortest}{The exact conditional-family optimum from the source}
The fixed-placement interval \(I_*^{\rm T}\) minimizes false
coverage for every true \((s,\eta)\) and every \(t\ne s\),
with \(t>0\),
among all procedures whose acceptance section in
\(0<x_2<t-x_1\) has length \(\alpha(t-x_1)\) for almost every
\(0<x_1<t\).
This includes noninterval sections and the switched construction.
It is the source's conditional-similarity class, not a claim that
every possible similar region belongs to that class.

Put \(\rho_s=s^3\eta\). The achieved false-coverage probability is
\[
             \alpha z^2\left(1-\frac{\rho_s}{2}(1-z)\right),
\]
where
\[
 z=\begin{cases}
 t/s,&t\leq s,\\
 [1-(1-\alpha)t/s]/\alpha,&s<t<s/(1-\alpha),\\
 0,&t\geq s/(1-\alpha).
 \end{cases}
\]
It equals \(\alpha\) at \(t=s\) and is less than \(\alpha\)
otherwise.
\end{hbtheorem}

\begin{hbtheorem}{thm:p15-neyman-regular-obstruction}{proved=annex:p15.neyman.regular.obstruction}{The regular two-sided obstruction with explicit analytic premises}
Let \(p_s\) be probability densities on \(\mathbb R^d\),
\(d\geq1\), with \(s\) varying on an open interval about \(t\).
Assume \(s\mapsto p_s\) is differentiable at \(t\) in \(L^1\),
with derivative \(g\) satisfying \(\int|g|>0\).
There is no measurable region \(A\) with \(P_t(A)=\alpha\)
which minimizes \(P_s(A)\) among all regions of that same
\(P_t\)-probability for every \(s\ne t\) in a two-sided
neighborhood of \(t\).
This gives the regular obstruction underlying the source's
Proposition X. Nonzero derivative is required at the null
parameter itself, and the differentiability permits integration
over the competing regions.
\end{hbtheorem}

\begin{hbtheorem}{thm:p15-neyman-normal-one-sided}{proved=annex:p15.neyman.normal.one.sided}{A complete normal example of one-sided optimality and failure of a shortest two-sided system}
Let \(X\sim\operatorname{normal}(\mu,\sigma^2)\), with known
\(\sigma>0\), and let \(z_\alpha\) be the standard normal
\(\alpha\)-quantile, so \(\Phi(z_\alpha)=\alpha\) for the
standard normal cumulative function \(\Phi\). The one-sided intervals
\[
             [X-\sigma z_\alpha,\infty),\qquad
             (-\infty,X+\sigma z_\alpha]
\]
each have coverage \(\alpha\) and satisfy the corresponding
best-one-sided criterion.
No exact-\(\alpha\) interval procedure is a shortest system
against all two-sided false targets in this experiment.
This supplies an explicit regular example; it does not replace
the source's unknown-variance Student problem.
\end{hbtheorem}

\begin{hbdefinition}{def:p15-neyman-unbiased-interval}{The source's global unbiasedness and local shortness conditions}
An exact-\(\alpha\) system is unbiased in the source's sense if,
for every fixed candidate \(t\),
\[
                   P_{s,\eta}\{t\in I(X)\}\leq\alpha
                              \quad\hbox{for every true }s,\eta,
\]
with equality at \(s=t\).
This is a property of coverage probabilities, not mean-unbiasedness
of an endpoint.
Suppose \(P_{s,\eta}\) has density \(p_{s,\eta}\) relative to
a specified common measure \(\nu\). Where two differentiations
under the acceptance integral are justified, the source's local shortness comparison minimizes
\[
             \left.\frac{\partial^2}{\partial s^2}
                            P_{s,\eta}(A_I(t))\right|_{s=t}
\]
among the declared exact-coverage unbiased class, with first
derivative zero, for every nuisance value.
\end{hbdefinition}

\begin{hbtheorem}{thm:p15-neyman-unbiased-derivatives}{proved=annex:p15.neyman.unbiased.derivatives}{What the derivative conditions establish}
For an interior target, global unbiasedness and the stated
twice differentiability imply
\[
 \int_{A_I(t)}\left.\partial_s p_{s,\eta}\right|_{s=t}\,d\nu=0,\qquad
 \int_{A_I(t)}\left.\partial_s^2 p_{s,\eta}\right|_{s=t}\,d\nu\leq0.
\]
Minimizing the second integral is the original local curvature
comparison. These derivative conditions by themselves do not
imply global unbiasedness.
For example, in the binary experiment with
\[
 P_s(X=1)=\tfrac12-s^2+4s^4,\qquad |s|<3/5,
\]
the acceptance region \(\{1\}\) has probability \(1/2\) at zero,
first derivative zero and second derivative \(-2\) there,
but has probability greater than \(1/2\) for some other
allowed parameters.
The two explicitly optimal uniform and triangular systems above
do satisfy the global unbiasedness inequalities.
\end{hbtheorem}

The unknown-variance normal construction in Neyman's original Example III
is supplied in \cref{sec:p15-student-inference}, including its Student
pivot, exact coverage and both original comparison criteria. Its global
unbiased optimality is established by a separate full argument.
